\subsection{Complete-request benefit}
\label{sec:request}


\begin{table}[t]
\centering\small
\begin{tabular}{lrr}
\toprule
Executor & 32 steps (ms) & 64 steps (ms) \\
\midrule
Native shape cache & \res{Native32} & \res{Native64} \\
FuseOS persistent & \res{Persistent32} & \res{Persistent64} \\
CLC+PDL shape cache & \res{CLC32} & \res{CLC64} \\
CLC+PDL fixed topology & \res{Fixed32} & \res{Fixed64} \\
\bottomrule
\end{tabular}
\caption{Qualified complete requests: median of three process medians,
five requests per cell. Each interval includes fresh prefill and output-driven
search. All methods reproduce the native final paths.}
\label{tab:requests}
\end{table}

The FuseOS executor retains an application-level benefit after the attention repair.
For 32 search steps, the complete request takes \res{Persistent32}\,ms,
compared with \res{Native32}\,ms for native shape-cached execution
(\res{RequestGain32}$\times$). For 64 steps, the corresponding times are
\res{Persistent64}\,ms and \res{Native64}\,ms
(\res{RequestGain64}$\times$). The complete-request comparison includes the native
GPU-length graph, which keeps graph count bounded by frontier width rather
than sequence depth. Thus the comparison is not against repeated Python
launches or compulsory graph reconstruction.

All 28 decoder layers and the vocabulary head execute in these intervals.
The searches expand 128 and 251 branch-tokens respectively; the latter
changes the live frontier between three and four branches. Both use model
outputs for parent/token selection, and the corrected custom paths reproduce
the final native candidate paths. The maximum request-gate NRMSE is
\res{MaxRequestNRMSE}, below the unchanged threshold.

\paragraph{What the positive result establishes.}
The complete dynamic executor is useful for this narrow-frontier workload.
The result does not isolate persistence from fused arithmetic. With the same
fused handlers, shape-cached CLC+PDL completes the two requests in
\res{CLC32}\,ms and \res{CLC64}\,ms, faster than the persistent executor.
The fixed-topology CLC+PDL control takes \res{Fixed32}\,ms and
\res{Fixed64}\,ms. This comparison distinguishes the cost of the execution
protocol from the benefit of the fused implementation. We retain the full
run ranges, including prefill variation, instead of presenting the small
fixed-topology differences as a universal persistent advantage.

\subsection{Dependence on workload geometry}
\label{sec:geometry}


We next hold the search trajectory constant to remove work differences from
the backend comparison. Three configurations vary ancestor length $P$,
maximum frontier width $B$, and horizon $H$: $(128,4,32)$, $(512,16,16)$, and
$(1024,8,32)$. Each uses a native-generated, output-dependent trace. All
qualified executors perform the same selected histories and execute the
selection policy during timing.

The equal-work replay shows that the narrow-frontier benefit does not
extend uniformly to wider or longer-prefix cases. Persistence improves on
the matched tile-grid and tile+PDL implementations, but native cached
execution becomes competitive or faster as workload geometry changes.
Shape-cached CLC+PDL provides another strong control. Since prefix, width,
horizon, and trajectory all change between cases, this comparison defines an
observed applicability boundary rather than identifying a single causal
shape parameter. The maximum post-fix replay NRMSE is
\res{MaxReplayNRMSE}; timing includes no gate-failing configuration.

